\chapter{Cultures et variétés}
\label{chap:cultures}

\section{Vue d'ensemble}

Une \emph{culture} (\emph{crop}) regroupe les variétés d'une même espèce botanique : « Tomate », « Carotte », « Pommier ». Chaque culture appartient à une famille botanique (Solanacées, Apiacées, Rosacées…) et à une strate de végétation.

Une \emph{variété} (\emph{variety}) est une déclinaison d'une culture : « Tomate Coeur-de-bœuf », « Tomate Marmande », « Pommier Reine des Reinettes ». Les variétés portent les caractéristiques agronomiques précises (durée jusqu'à transplant, durée jusqu'à maturité, jours-de-l'année de récolte pour les pluriannuelles).

\section{L'écran Cultures}

\begin{todobox}
À détailler : layout maître-détail (liste cultures à gauche, variétés à droite), création d'une culture, choix annuel/pluriannuel.
\end{todobox}

\section{Créer une variété}

\begin{todobox}
À détailler : champs spécifiques annuelles (DTT, DTM, fenêtre de récolte) vs pluriannuelles (DOY débourrement, floraison, début/fin récolte, rendement attendu).
\end{todobox}
